\section{Conclusion}
\label{sec:conclusion}

Supplier competition improves procurement by lowering expected production cost. This paper shows that the same competition can worsen a different margin: decentralized coordination over which region undertakes a complementary industrial-policy role.

The mechanism relies on incidence. More suppliers raise expected real surplus from a complementary project but, under standard DRHR procurement conditions, reduce winning supplier information rent. A coordinated policymaker responds to the first effect, while a supplier-side regional government responds to the locally captured second effect. Consequently, competition lowers the coordinated threshold for regional differentiation but raises the decentralized threshold. The parameter region with unique inefficient duplication expands.

The stronger result is an equilibrium reversal. If differentiated production remains socially valuable when the supplier market becomes very competitive, then sufficiently strong supplier competition eliminates differentiated Nash equilibria and leaves duplicated local priorities as the unique decentralized outcome. The market becomes more efficient conditional on production while decentralized policy selection moves away from the coordinated allocation.

The auction-theoretic rent-dissipation result is not new, and neither is the generic observation that incomplete local capture can distort decentralized decisions. The contribution is their interaction across strategic layers: a familiar competition effect in the supplier market can change the equilibrium set of regional industrial policy. This mechanism suggests that procurement reform and industrial-policy coordination should be evaluated jointly when complementary production roles span jurisdictions.
